\documentclass[../../../include/open-logic-section]{subfiles}

\begin{document}

\olfileid{sth}{replacement}{limofsize}
\olsection{పరిమాణ పరిమితి}

ప్రతిస్థాపనకు ``అంతర్గత'' సమర్థన ఇచ్చే అత్యంత సాధారణ ప్రయత్నం బహుశా ఈ భావన ద్వారా వస్తుంది:
\begin{enumerate}
	\item[] \limofsize. వస్తువులు మరీ ఎక్కువగా లేకపోతే, ఏ వస్తువులైనా ఒక సమితిని ఏర్పరుస్తాయి.
\end{enumerate}
ఈ సూత్రం వెంటనే ప్రతిస్థాపనను సమర్థిస్తుంది. ఎందుకంటే ప్రతిస్థాపనతో ఏర్పడిన సమితి, దాన్ని ఏర్పరచడానికి వాడిన సమితికంటే పెద్దది కాలేదు. ఖచ్చితంగా చెప్పాలంటే: ప్రతిస్థాపనతో $\funimage{\tau}{A} = \Setabs{\tau(x)}{x \in A}$ సమితిని ఏర్పరచారనుకుందాం; అప్పుడు $\cardle{\funimage{\tau}{A}}{A}$; కాబట్టి $A$లోని \tetoken{మూలకాలు}{!!{element}s} సమితిని ఏర్పరచలేనంత ఎక్కువ కాకపోతే, వాటి బింబాలు కూడా $\funimage{\tau}{A}$ను ఏర్పరచలేనంత ఎక్కువ కావు.

ప్రతిస్థాపనను సమర్థించడానికి \limofsize{}ను వాడడంలో స్పష్టమైన కష్టం: ఇంతవరకు మనం \limofsize{} వంటి ఏ సూత్రాన్నీ \emph{ప్రవేశపెట్టలేదు}. అంతేకాక, \crefrange{sth:story::chap}{sth:z::chap}లో సమితి సంచిత-పునరావర్తన భావన కథనం చెప్పినప్పుడు, ఏదీ \limofsize{} వైపు \emph{సూచించలేదు}. అందుకే బూలాస్ ఒక సందర్భంలో ఇలా రాశారు: ``సమితి సిద్ధాంతం `వెనుక' కనీసం రెండు ఆలోచనలు ఉన్నాయని బహుశా తేల్చవచ్చు'' \citeyearpar[p.~19]{Boolos1989}. ఒక వైపు, సంచిత-పునరావర్తన సమితి భావనకు సంబంధించిన ఆలోచనలు~$\Z$ను సమర్థించడానికి ఉద్దేశించబడ్డాయి. మరో వైపు, \limofsize{} ప్రతిస్థాపనను సమర్థించడానికి ఉద్దేశించబడింది.

అయితే సమస్య ఇంతవరకు \limofsize{} గురించి మనం \emph{మౌనంగా} ఉండటమే కాదు. ఇప్పుడే చెప్పిన రూపంలో \limofsize{}, సంచిత-పునరావర్తన సమితి భావనతో సరిగా పొసగనట్టు కనిపించడం అసలు సమస్య. ఎందుకంటే సమితులు \emph{దశల్లో} ఏర్పడతాయనే భావనను అది అసలు ప్రస్తావించదు.

ఈ ``రెండు ఆలోచనల'' చరిత్రను చూస్తే ఆశ్చర్యం లేదు (అంటే సంచిత-పునరావర్తన సమితి భావన, \limofsize). చివరకు ఇవి సమితి సిద్ధాంత వైరుధ్యాలను అడ్డుకునే రెండు భిన్నమైన కార్యక్రమాలు. సంచిత-పునరావర్తన భావన రస్సెల్ వైరుధ్యాన్ని స్థూలంగా ఇలా అడ్డుకుంటుంది: \emph{రస్సెల్ సమితి ఉంటుందని మనం ఎప్పుడూ ఆశించకూడదు, ఎందుకంటే అది ఏ దశలోనూ ``ఏర్పడదు''}. దీనికి భిన్నంగా, \limofsize{} రస్సెల్ సమితిని స్థూలంగా ఇలా నిరాకరిస్తుంది: \emph{రస్సెల్ సమితి ఉంటుందని మనం ఎప్పుడూ ఆశించకూడదు, ఎందుకంటే అది మరీ పెద్దదిగా ఉంటుంది}.

ఇలా చెప్పాక నేరుగా అనుకుందాం: వైరుధ్యాలకు ప్రతిస్పందనగా పరిగణిస్తే, \limofsize{}కు \emph{చాలా} ఎక్కువ సమర్థన అవసరం. ఉదాహరణకు రస్సెల్ వైరుధ్యం యొక్క ఈ రూపాన్ని చూడండి: \emph{తమను తాము వాసన చూడని పగ్ కుక్కలను మాత్రమే వాసన చూసే పగ్ కుక్క లేదు} (\olref[sth][story][rus]{sec} చూడండి). ``అలాంటి పగ్ కుక్క ఎందుకు లేదు?'' అని అడిగితే, అది మరీ ఎక్కువ పగ్ కుక్కలను వాసన చూడాల్సి వస్తుందని చెప్పడం మంచి సమాధానం కాదు. మరి రస్సెల్ సమితి లేకపోవడానికి అది ఉనికిలో ఉండలేనంత ``పెద్దదిగా'' ఉండాల్సి వస్తుందని చెప్పడం మంచి సహజ వివరణ ఎందుకవుతుంది?

సంక్షిప్తంగా, \limofsize{}కు ``సహజ'' ప్రేరణ ఏమిటో కొంత అయోమయంగా అనిపిస్తే అది అర్థం చేసుకోదగినదే.

\end{document}
